\lysh{20865_SOLUTION}{1}
\textbf{ΛΥΣΗ}

α) Η εξίσωση της έλλειψης $C$ με εστίες τα σημεία $E'(-\gamma, 0), E(\gamma, 0)$, μήκος μεγάλου άξονα $2\alpha$ και μήκος μικρού άξονα $2\beta$, είναι
\[ \frac{x^2}{\alpha^2} + \frac{y^2}{\beta^2} = 1, \text{ όπου } \beta = \sqrt{\alpha^2 - \gamma^2} \]
Η εξίσωση της έλλειψης $C_1: x^2 + 4y^2 = 4$ γίνεται ισοδύναμα $\frac{x^2}{4} + y^2 = 1$. Είναι λοιπόν:
\[ \alpha^2 = 4 \underset{\alpha>0}{\iff} \alpha = 2 , \beta^2 = 1 \underset{\beta>0}{\iff} \beta = 1 \text{ και } 1 = \sqrt{4 - \gamma^2} \iff \gamma^2 = 3 \underset{\gamma>0}{\iff} \gamma = \sqrt{3} . \]
Επομένως, είναι:
\[ 2\alpha = 4, 2\beta = 2, E'(-\sqrt{3}, 0), E(\sqrt{3}, 0) . \]
Για την έλλειψη $C_2: \frac{x^2}{16} + \frac{y^2}{4} = 1$ είναι αντίστοιχα:
\[ \alpha^2 = 16 \underset{\alpha>0}{\iff} \alpha = 4, \beta^2 = 4 \underset{\beta>0}{\iff} \beta = 2 \text{ και } 2 = \sqrt{16 - \gamma^2} \iff \gamma^2 = 12 \underset{\gamma>0}{\iff} \gamma = 2\sqrt{3} . \]
Επομένως, είναι:
\[ 2\alpha = 8, 2\beta = 4, E'(-2\sqrt{3}, 0), E(2\sqrt{3}, 0) . \]

β) Δύο ελλείψεις όταν έχουν την ίδια εκκεντρότητα $\varepsilon$ λέγονται όμοιες, όπου $\varepsilon = \frac{\gamma}{\alpha}$.

Η έλλειψη $C_1$ έχει εκκεντρότητα $\varepsilon_1 = \frac{\sqrt{3}}{2}$.

Η έλλειψη $C_2$ έχει εκκεντρότητα $\varepsilon_2 = \frac{2\sqrt{3}}{4} = \frac{\sqrt{3}}{2}$.

Επομένως, ο ισχυρισμός είναι αληθής.
\end{lysh}